% Generated draft: requires mathematical and source review.
\documentclass[11pt]{article}
\usepackage[margin=25mm]{geometry}
\usepackage{fontspec}
\setmainfont{Latin Modern Roman}
\usepackage{amsmath,amssymb,mathtools,unicode-math}
\setmathfont{Latin Modern Math}
\usepackage{enumitem,needspace,xurl,hyperref}
\setlength{\parindent}{0pt}
\setlength{\parskip}{5pt}
\setlength{\emergencystretch}{3em}
\newcommand{\N}{\mathbb{N}}
\newcommand{\Z}{\mathbb{Z}}
\newcommand{\Q}{\mathbb{Q}}
\newcommand{\R}{\mathbb{R}}
\newcommand{\C}{\mathbb{C}}
\newcommand{\F}{\mathbb{F}}
\newcommand{\A}{\mathbb{A}}
\newcommand{\PP}{\mathbb{P}}
\newcommand{\E}{\mathbb{E}}
\newcommand{\eps}{\varepsilon}
\newcommand{\dd}{\,\mathrm{d}}
\begin{document}
\section*{UnsolvedMath ID: 30002033}
\textbf{Source problem code:} catalogue entry
\subsection{Definitions and mathematical statement}
The supplied record combines two questions. The arithmetic-progression question admits the following precise formulation under an explicit counting convention; the graph-structure question remains underspecified.

For each integer \(n\geq 1\), let \([n]=\{1,\ldots,n\}\). A two-colouring is a map \(\chi:[n]\to\{0,1\}\). Adopt the convention that a three-term arithmetic progression is a triple \((a,a+d,a+2d)\) with integers \(a\geq 1\), \(d\geq 1\), and \(a+2d\leq n\), counted once in increasing order. Define
\[
M_n(\chi)=\left|\left\{(a,d)\in\Z^2:a\geq1,\ d\geq1,\ a+2d\leq n,\ \chi(a)=\chi(a+d)=\chi(a+2d)\right\}\right|
\]
and
\[
m(n)=\min_{\chi:[n]\to\{0,1\}}M_n(\chi).
\]
The arithmetic-progression problem is to determine the asymptotically minimal count \(m(n)\) as \(n\to\infty\), including its leading coefficient if it has the form
\[
m(n)=c n^2+o(n^2).
\]
This formulation does not assert that such a constant \(c\) exists. The source does not specify its counting or density normalization. The convention above excludes constant progressions and identifies the two orientations of a progression; it must be checked against the original report before numerical constants are compared.

The other source question asks whether an analogue of the ``quadratic-structure phenomenon'' for two-colourings with fewer than the random density of monochromatic \(K_4\)'s holds for graph colourings. Here \(K_4\) denotes the complete graph on four vertices. The supplied record does not define the quoted phenomenon, the objects coloured in its premise, the graph-colouring model in its conclusion, the random comparison model, or the structural conclusion and its quantifiers. Consequently, this part cannot be converted into a precise mathematical assertion from the supplied evidence alone; no structural theorem or conjectural implication is being supplied here.
\subsection{Short English statement}
Determine the smallest possible number of monochromatic three-term arithmetic progressions in a two-colouring of \(\{1,\ldots,n\}\), asymptotically as \(n\to\infty\), using the stated counting convention. The accompanying question about quadratic structure and monochromatic \(K_4\)'s needs additional definitions from its source. Both parts remain statement-only drafts requiring human mathematical review.
\subsection{Sources}
\begin{itemize}
\item Supplied local UnsolvedMath registry record: permanent ID \(30002033\), title ``Monochromatic Structures below Random Density'', source problem code \texttt{catalogue entry}, collection ``Oberwolfach Reports Open Problems''. No individual proposer is supplied; the proposed year is recorded as 2012.
\item Supplied report citation: \emph{Mini-Workshop: Hypergraph Tur\'an Problem (2013)}, Oberwolfach Reports, pp.~1149--1182, DOI 10.4171/owr/2012/19. \url{https://doi.org/10.4171/owr/2012/19}. The record describes Oberwolfach Reports as documenting mathematical workshops and their open-problem discussions. Its citation year 2013 differs from its proposed-year and year-tag value 2012; this discrepancy has not been resolved.
\item Supplied literature reference: Parrilo--Robertson--Saracino. \url{https://doi.org/10.1016/j.jcta.2007.03.006}.
\item Supplied literature reference: Greenwood--Kariv--Williams. \url{https://arxiv.org/abs/2301.00336}.
\item Supplied literature reference: Oberwolfach report. \url{https://ems.press/content/serial-article-files/46390?nt=1}.
\item Literature assessment dated exactly \texttt{2026-08-21}: supplied dated curation, not independently verified. The supplied status is ``open''. The assessment states that the asymptotic minimum density of monochromatic three-term progressions in two-colourings is undetermined and that known bounds and restricted-class optimality results leave the general problem open. The background describes this as dated literature triage derived from the supplied Oberwolfach curation and requests independent verification of time-sensitive status and cited claims.
\item The registry labels the statement \texttt{corrected\_verified}, with a supplied note reporting manual semantic, typographic, and LaTeX review of a meaning-preserving cleanup. This is a source-record qualification, not independent verification of the formulation, literature, or mathematical status.
\end{itemize}
\end{document}
